\section{Global momentum uniqueness and the stable optical weight}
\label{sec:v99-uniqueness}
The multiplicity in \eqref{eq:v98-weighted-multiplicity} can be
estimated dynamically. Two separate facts are needed: a global
uniqueness assertion for each lifted itinerary, and a comparison of
the optical weight with a genuine stable Jacobian. Neither follows
merely from the nonzero local endpoint Hessian.

\begin{theorem}[One physical orbit per itinerary and momentum pair]
\label{thm:v99-unique}
Fix a finite lifted sequence of circular scatterers and two endpoint
tangential momenta in $(-1,1)$. There is at most one regular physical
trajectory with that itinerary and those endpoint momenta. Internal
incidences need not be bounded below uniformly. The assertion concerns
physical contact points, not their overlapping angular coordinates.
\end{theorem}
\begin{proof}
Write $Q=(Q_0,\ldots,Q_m)$ for the contact positions and
\[
                 S(Q)=\sum_{j=0}^{m-1}|Q_{j+1}-Q_j|.
\]
This is a convex function of the Cartesian contact variables. It is
differentiable at every configuration with positive flight lengths.
For two such configurations $Q,\widehat Q$, monotonicity of its
subgradient gives
\begin{equation}\label{eq:v99-cartesian-monotonicity}
 \sum_{j=0}^m
 (\nabla_j S(Q)-\nabla_j S(\widehat Q))\cdot(Q_j-\widehat Q_j)
                                                   \ge0.
\end{equation}
Let $n_j$ be the outward unit normal of the $j$th circular scatterer
and let $t_j$ be its fixed oriented tangent. Reflection and first
variation give
\[
 \nabla_0S=-c_0n_0-p_0t_0,\quad
 \nabla_jS=-2c_jn_j\ (0<j<m),\quad
 \nabla_mS=-c_mn_m+p_mt_m.
\]
At the two trajectories $c_0=\widehat c_0>0$ and
$c_m=\widehat c_m>0$, since their endpoint momenta agree. The
scatterer centers agree at every index, so
$Q_j-\widehat Q_j=R(n_j-\widehat n_j)$. Rotation by a right angle
is skew symmetric; consequently
$(t_j-\widehat t_j)\cdot(n_j-\widehat n_j)=0$.
The left side of \eqref{eq:v99-cartesian-monotonicity} therefore equals
\begin{align*}
 &-Rc_0|n_0-\widehat n_0|^2-Rc_m|n_m-\widehat n_m|^2\\
 &\hspace{12mm}-R\sum_{j=1}^{m-1}
                  (c_j+\widehat c_j)|n_j-\widehat n_j|^2.
\end{align*}
Every coefficient is strictly positive. Thus all contact positions
coincide. The flight directions and physical state then coincide as
well. Repeated visits to a scatterer cause no change: the contact
variables at different times are distinct Cartesian variables. No
convexity of a physical word domain, existence of an interpolating
physical trajectory, or global angular lift is needed.
\end{proof}

For the Jacobian comparison use standard billiard coordinates
$(s,\varphi)$, where $s$ is arclength and $|p|=|\sin\varphi|$.
Signs are fixed consistently on each oriented boundary component.
Let $M'$ be the full-measure set regular in both infinite time
directions. Write $J^u_RT_R^j$ and $J^s_RT_R^j$ for the positive
arclength Jacobians on its invariant unstable and stable lines.
The invariant cones have slopes in fixed compact positive and
negative intervals, respectively, with a uniform angle between them
\cite[Section 1.1]{BDPressure}.

\begin{lemma}[Optical weight and removal of one incidence]
\label{lem:v99-stable-weight}
There are uniform geometric constants $C_0,C_1$ such that on $M'$,
for every $m\ge1$,
\begin{equation}\label{eq:v99-stable-weight}
 C_0^{-1}c_m J^s_RT_R^m(x)\le\omega(x)
                         \le C_0c_m J^s_RT_R^m(x),\qquad
 \frac{J^s_RT_R(x)}{c(x)}\le C_1.
\end{equation}
In particular, for an internal mark $1\le i<m$,
\begin{equation}\label{eq:v99-remove-mark}
 \frac{\omega(x)}{c_i(x)}
 \le C_0C_1 J^s_RT_R^i(x)
                    J^s_RT_R^{m-i-1}(T_R^{i+1}x).
\end{equation}
The stable factor at the marked time is removed, not evaluated with
a singular inverse-incidence multiplier left over.
\end{lemma}
\begin{proof}
Up to consistent orientation signs, one flight has derivative
\[
 \frac1{c'}
 \begin{pmatrix}
 \tau\varkappa+c&\tau\\
 \tau\varkappa\varkappa'+\varkappa c'+\varkappa'c&
                         \tau\varkappa'+c'
 \end{pmatrix},\qquad \varkappa=\varkappa'=R^{-1}.
\]
Its determinant is $c/c'$. Every entry is positive. The ratio of
its first column to its second, in either row, lies between
$\varkappa_{\min}$ and
$\varkappa_{\max}+\tau_-^{-1}$. Left multiplication by a positive
matrix preserves these componentwise inequalities. If $b_m$ is the
upper-right entry of the $m$-flight product, the image of the vertical
vector has slope in the unstable cone. Comparing it with a unit
invariant unstable vector at the source consequently gives
\[
           C^{-1}b_m\le J^u_RT_R^m(x)\le Cb_m.
\]
Here the constants do not accumulate with $m$: both the initial
column comparison and the final slope comparison use fixed cones.
Changing from $\varphi_0$ to $p_0$ and from $s_m$ to $\alpha_m$
gives $\omega=Rc_0/b_m$.

Put $e(x)=\sin\angle(E^s(x),E^u(x))$. Conservation of collision
area in these coordinates is the exact identity
\[
 J^u_RT_R^m(x)J^s_RT_R^m(x)
                    \frac{e(T_R^mx)}{e(x)}=\frac{c_0}{c_m}.
\]
The function $e$ is bounded above and below by positive cone
constants. Combining the last three formulas proves the two-sided
comparison in \eqref{eq:v99-stable-weight}.
For one flight, if $v$ is the initial unstable slope, then
\[
 J^u_RT_R(x)=
 \frac{\tau\varkappa+c+\tau v}{c'}
                  \frac{\sqrt{1+(v')^2}}{\sqrt{1+v^2}}.
\]
The same area identity shows that $J^s_RT_R/c$ is bounded by a
fixed constant divided by $\tau\varkappa+c+\tau v$. The latter is
bounded below by a positive geometric constant. This proves the
one-step bound. Factor the stable cocycle at time $i$ and use
$c_m\le1$ to obtain \eqref{eq:v99-remove-mark}.
\end{proof}

The area formula in module 235 can be restricted to $M'$ without
altering its almost-everywhere momentum density. The omitted set has
zero collision area. On countably many regular endpoint charts $G$
is a local diffeomorphism, so its image has zero two-dimensional
Lebesgue measure. Thus the infinite-time invariant lines used above
do not add a restriction to the original finite-record source.
